% !TEX root = ../../main.tex

\subsection{Notation}

We use one notation layer throughout the method to keep the statistical and
algorithmic steps aligned. The working sample representation is a matrix $X$
equipped with an explicit feature-space contract, the hierarchy is a rooted
binary tree $\mathcal{T} = (V,E)$, and the main inferential objects are
node-level distributions $\theta_u$ estimated from descendant leaves. Edge-stage
quantities compare a child $c$ to its parent $u$, sibling-stage quantities
compare $l(u)$ and $r(u)$ under the same parent, and traversal quantities
convert the resulting significance indicators into the final partition.
\Cref{tab:notation} summarizes the symbols that carry the rest of the method.

{\small
\begin{longtable}{@{}p{0.24\textwidth}p{0.70\textwidth}@{}}
\caption{Core notation used throughout the KL-TE method description.}
\label{tab:notation}\\
\toprule
Symbol & Meaning \\
\midrule
\endfirsthead
\toprule
Symbol & Meaning \\
\midrule
\endhead
\bottomrule
\endfoot
$n$ & Number of samples. \\
$p$ & Number of raw working coordinates before block-specific contrast charts. \\
$X \in \mathbb{R}^{n \times p}$ & Numeric working sample-feature matrix with a Bernoulli, categorical, continuous, or mixed feature-space contract. \\
$X_i \in \mathbb{R}^p$ & Row vector for sample $i$. \\
$\mathcal{T} = (V,E)$ & Rooted binary hierarchy built from the samples. \\
$u, v, a$ & Generic tree nodes. \\
$L_i$ & Leaf node corresponding to observed sample $i$. \\
$R$ & Root node whose descendant set contains all samples. \\
$c$ & Generic child node in a parent-child edge comparison. \\
$l(u), r(u)$ & Left and right child of parent node $u$. \\
$L(u)$ & Descendant leaf set of node $u$. \\
$n_u = |L(u)|$ & Number of descendant leaves under node $u$. \\
$e = (u,c)$ & Directed tree edge from parent $u$ to child $c$, representing immediate subtree containment. \\
$\theta_u$ & Empirical node distribution at node $u$ in the active feature-space coordinates. \\
$\beta_u=n_{l(u)}/n_u$ & Left-child barycentric weight at binary parent $u$. \\
$b_u=\min(\beta_u,1-\beta_u)$ & Barycentric balance of the two children at parent $u$. \\
$s_u^{\mathrm{samp}}=1/n_{l(u)}+1/n_{r(u)}$ & Two-sample sampling variance scale for the sibling contrast at parent $u$. \\
$\ell_u^{\mathrm{bar}}$ & Barycentric leverage, defined as $\log\{\max(\beta_u,1-\beta_u)/\min(\beta_u,1-\beta_u)\}$. \\
$H_{0,j}^{\mathrm{edge}}(u,c)$ & Coordinatewise child-parent edge null for feature $j$. \\
$z_{u \to c}^{\mathrm{edge}}$ & Standardized child-parent edge vector. \\
$T_{u \to c}^{\mathrm{edge}}$ & Projected edge statistic for edge $(u,c)$. \\
$p_{u \to c}^{\mathrm{edge}}$ & Raw edge p-value. \\
$p_{u \to c}^{\mathrm{edge,BH}}$ & Tree-BH-adjusted edge p-value used for edge decisions and sibling calibration weights. \\
$S^{\mathrm{edge}}(c)$ & Edge significance indicator after multiplicity correction. \\
$H_0^{\mathrm{sib}}(u)$ & Sibling null at parent $u$. \\
$\bar\theta_{u,j}$ & Pooled sibling proportion at coordinate $j$. \\
$z_u^{\mathrm{sib}}$ & Standardized sibling contrast vector at parent $u$. \\
$k_u^{\mathrm{edge}}, k_u^{\mathrm{sib}}$ & Edge-stage and sibling-stage projection dimensions. \\
$T_u^{\mathrm{raw}}$ & Raw projected sibling statistic. \\
$a_u$ & Projected quadratic reference scale; $a_u=1$ for the current orthonormal PCA basis. \\
$\nu_u$ & Sibling reference degrees of freedom; $\nu_u=k_u^{\mathrm{sib}}$ for the current orthonormal PCA basis. \\
$x_q=(\log s_q,\log n_q)$ & Sibling calibration context used by the empirical-null inflation estimator. \\
$s_q$ & Resolved sibling projection dimension for sibling record $q$. \\
$\pi_q$ & Empirical sibling null weight used in inflation estimation. \\
$w_q(u)$ & Context weight assigned to calibration record $q$ when estimating inflation for sibling parent $u$. \\
$h$ & Null-weighted bandwidth vector used on active sibling calibration context axes. \\
$n_{\mathrm{cal}}, n_{\mathrm{eff}}$ & Number of contributing sibling calibration records and their effective sample size. \\
$\hat c_{f_u}(x_u)$ & Context-weighted empirical-null inflation estimate for sibling parent $u$ in feature family $f_u$. \\
$S^{\mathrm{sib}}(u)$ & Sibling significance indicator after correction. \\
$\mathrm{Split}(u)$ & Indicator that traversal keeps the split at node $u$. \\
$D^{\mathrm{sib}}(u)$ & Indicator that some proper descendant of $u$ has a significant sibling split among nodes where the sibling test was run. \\
$\mathcal{C}$ & Final set of terminal cluster nodes returned by traversal. \\
\end{longtable}
}
